\section{22 September 2026: Broome and the next ten boundary reviews}

GO Broome's earlier building floor is zero. City Planning records complete
demolition in June 2019, before the November reference snapshot and March
2020 filing. Following Jacob's instruction, the manual source table explicitly
corrects the stale 4,600-square-foot archive value. Its 32,395 square feet of
ground and archived prior-use category remain unchanged.

The next ten parents were selected from the 116 unresolved weighting-sample
sites, prioritizing exact-99 constituents, then constituent count, parent units,
and parent ID. Five yield usable production corrections. Charles gains omitted
old lot 7, raising ground from 10,000 to 12,590 square feet and earlier floor
from zero to 1,100. Its older-proposal question remains open. Walton/Burnside's
recorded zoning diagram supports the complete earlier 23,000-square-foot lot;
its direct administrative floor is 21,000. Willets Phase I uses its recorded
135,609-square-foot development boundary, five contributing earlier parcels,
and zero earlier floor. Senior housing and the stadium occupy separate sites.

For 261/315 Grand Concourse, the DOF change names both DOB jobs and old lots
1/11/27. The recorded survey measures 37,661.3 square feet; their three earlier
floor records sum to 46,463. Noble/Oak uses the city's approximate 173,834
square feet and the complete 1,500 earlier floor. An independent reconstruction
of its recorded deed gives 173,827.90 square feet, with a 0.0082-foot closing
error. That agreement supports the development boundary within the larger
waterfront tax parcel. These corrections use frozen prefiling zoning and do
not introduce new floor-share imputations.

Four of the ten receive complete boundary decisions, reducing unresolved
parents from 116 to 112 among the same 864 sites. Every other pass/flag status
is unchanged. Charles, Broadway,
Livingston, Rockaway III/IV, Wallabout/Union, and 310/322 Grand Concourse retain
specific questions. The last two now have stronger ground evidence, but
allocation of earlier zoning or floor still requires work. Rockaway's newly
retrieved deed conveys condominium interests in land shared by several phases;
its ownership percentages do not measure each phase's physical ground.

Livingston also raises a sample-timing concern. Its first DOB filing is in
September 2022, but a recorded declaration reproduces an April 2025 HPD
application for the same job with 99 units and 485-x Option B. This establishes
post-policy treatment sought for the current design; it does not establish the
original 2022 count or amendment date. The current cohort assignment is preserved
in this boundary batch. Its interpretation as an untreated historical 99-unit
choice requires a dated filing-history review before estimation.

Exactly six parent rows change, including Broome. Every constituent record,
parent housing count, filing date, and membership is unchanged. The full panels
retain 1,811 historical parents and 907 post-policy parents; the weighting
sample retains 554 and 310. Under the existing calibration, the weighted
historical exact-99 share changes from 1.3535331 to 1.3537427 percent and mean
parent units from 151.4686 to 151.3559. Maximum moment error is
$6.50\times10^{-10}$. The six previously approved floor-estimate flags remain.
These small descriptive changes do not settle the remaining allocation or
timing questions for structural estimation.

\textit{Sources and reproduction.} The audit's
\texttt{next\_ten\_boundaries\_2026-09-22.md} records all ten decisions, original
document pages, and unresolved evidence. Unchanged ACRIS exports and search
extracts are preserved; nine public documents have concrete checksum-pinned
download recipes. Production reads the manual table and archived parcels;
22v1 is added as an explicit input for Willets. Root \texttt{make} rebuilds
the panels and exhibits, and \texttt{make logbook} refreshes the audits,
calibration, sensitivities, and this record.
